
\documentclass{article}
\usepackage{colortbl}
\usepackage{makecell}
\usepackage{multirow}
\usepackage{supertabular}

\begin{document}

\newcounter{utterance}

\twocolumn

{ \footnotesize  \setcounter{utterance}{1}
\setlength{\tabcolsep}{0pt}
\begin{supertabular}{c@{$\;$}|p{.15\linewidth}@{}p{.15\linewidth}p{.15\linewidth}p{.15\linewidth}p{.15\linewidth}p{.15\linewidth}}

    \# & $\;$A & \multicolumn{4}{c}{Game Master} & $\;\:$B\\
    \hline 

    \theutterance \stepcounter{utterance}  

    & & \multicolumn{4}{p{0.6\linewidth}}{\cellcolor[rgb]{0.9,0.9,0.9}{%
	\makecell[{{p{\linewidth}}}]{% 
	  \tt {\tiny [A$\langle$GM]}  
	 State giocando a un gioco di negoziazione in cui dovete accordarvi su come dividere una serie di oggetti tra voi e un altro giocatore.\\ \tt \\ \tt Regole:\\ \tt (a) A Lei e all'altro giocatore viene dato un insieme di oggetti. Ognuno di voi riceve anche una funzione di valore segreta, che rappresenta il valore di ciascun tipo di oggetto.\\ \tt (b) Si scambiano messaggi con l'altro giocatore per concordare chi si aggiudica gli oggetti. Potete inviare un massimo di 5 messaggi ciascuno, oppure terminare in anticipo facendo una proposta segreta in qualsiasi momento.\\ \tt (c) A ciascuno di voi viene chiesto di inviare una proposta segreta indicando gli oggetti che desiderate, formattata tra parentesi quadre come segue: "[Proposta: <numero> <nome oggetto>, <numero> <nome oggetto>, <...>]".\\ \tt (d) Se le vostre proposte sono complementari, cioè ci sono abbastanza oggetti per soddisfare entrambe le proposte, a ciascun giocatore viene assegnato un punteggio basato sulla somma dei valori degli oggetti ricevuti. In caso contrario, entrambi ricevono zero punti.\\ \tt (e) L'obiettivo è massimizzare il proprio punteggio.\\ \tt \\ \tt Cominciamo.\\ \tt \\ \tt L'insieme degli oggetti disponibili è:\\ \tt \\ \tt 2 fili, 3 bulloni, 1 quadro\\ \tt \\ \tt La funzione valore segreta è:\\ \tt \\ \tt filo: 0, bullone: 1, quadro: 7\\ \tt \\ \tt IMPORTANTE: i vostri messaggi, a meno che non si tratti di una proposta segreta, vengono trasmessi direttamente all'altro giocatore, quindi non includete alcuna risposta alle regole o testo di annuncio del vostro messaggio. Per fare una proposta segreta, utilizzate il formato indicato. Non utilizzare delle parentesi quadre quando si comunica all'altro giocatore, altrimenti verrà interpretata come una proposta segreta.\\ \tt  
	  } 
	   } 
	   } 
	 & \\ 
 

    \theutterance \stepcounter{utterance}  

    & \multicolumn{4}{p{0.6\linewidth}}{\cellcolor[rgb]{0.8,1,0.9}{%
	\makecell[{{p{\linewidth}}}]{% 
	  \tt {\tiny [A$\rangle$GM]}  
	 Ciao! Vedo che abbiamo 2 fili, 3 bulloni e 1 quadro da dividere tra noi. Per trovare un accordo equo, mi piacerebbe capire quali oggetti sono più importanti per te. Personalmente, il quadro mi interessa molto, mentre sono più flessibile sui fili e bulloni. E tu, hai delle preferenze particolari su questi oggetti? 
	  } 
	   } 
	   } 
	 & & \\ 
 

    \theutterance \stepcounter{utterance}  

    & & & \multicolumn{2}{p{0.3\linewidth}}{\cellcolor[rgb]{0.95,0.95,0.95}{%
	\makecell[{{p{\linewidth}}}]{% 
	  \tt {\tiny [GM$|$GM]}  
	 continue 
	  } 
	   } 
	   } 
	 & & \\ 
 

    \theutterance \stepcounter{utterance}  

    & & & \multicolumn{2}{p{0.3\linewidth}}{\cellcolor[rgb]{0.95,0.95,0.95}{%
	\makecell[{{p{\linewidth}}}]{% 
	  \tt {\tiny [GM$|$GM]}  
	 Ciao! Vedo che abbiamo 2 fili, 3 bulloni e 1 quadro da dividere tra noi. Per trovare un accordo equo, mi piacerebbe capire quali oggetti sono più importanti per te. Personalmente, il quadro mi interessa molto, mentre sono più flessibile sui fili e bulloni. E tu, hai delle preferenze particolari su questi oggetti? 
	  } 
	   } 
	   } 
	 & & \\ 
 

    \theutterance \stepcounter{utterance}  

    & & \multicolumn{4}{p{0.6\linewidth}}{\cellcolor[rgb]{0.9,0.9,0.9}{%
	\makecell[{{p{\linewidth}}}]{% 
	  \tt {\tiny [GM$\rangle$B]}  
	 State giocando a un gioco di negoziazione in cui dovete accordarvi su come dividere una serie di oggetti tra voi e un altro giocatore.\\ \tt \\ \tt Regole:\\ \tt (a) A Lei e all'altro giocatore viene dato un insieme di oggetti. Ognuno di voi riceve anche una funzione di valore segreta, che rappresenta il valore di ciascun tipo di oggetto.\\ \tt (b) Si scambiano messaggi con l'altro giocatore per concordare chi si aggiudica gli oggetti. Potete inviare un massimo di 5 messaggi ciascuno, oppure terminare in anticipo facendo una proposta segreta in qualsiasi momento.\\ \tt (c) A ciascuno di voi viene chiesto di inviare una proposta segreta indicando gli oggetti che desiderate, formattata tra parentesi quadre come segue: "[Proposta: <numero> <nome oggetto>, <numero> <nome oggetto>, <...>]".\\ \tt (d) Se le vostre proposte sono complementari, cioè ci sono abbastanza oggetti per soddisfare entrambe le proposte, a ciascun giocatore viene assegnato un punteggio basato sulla somma dei valori degli oggetti ricevuti. In caso contrario, entrambi ricevono zero punti.\\ \tt (e) L'obiettivo è massimizzare il proprio punteggio.\\ \tt \\ \tt Cominciamo.\\ \tt \\ \tt L'insieme degli oggetti disponibili è:\\ \tt \\ \tt 2 fili, 3 bulloni, 1 quadro\\ \tt \\ \tt La funzione valore segreta è:\\ \tt \\ \tt filo: 3, bullone: 0, quadro: 4\\ \tt \\ \tt IMPORTANTE: i vostri messaggi, a meno che non si tratti di una proposta segreta, vengono trasmessi direttamente all'altro giocatore, quindi non includete alcuna risposta alle regole o testo di annuncio del vostro messaggio. Per fare una proposta segreta, utilizzate il formato indicato. Non utilizzare delle parentesi quadre quando si comunica all'altro giocatore, altrimenti verrà interpretata come una proposta segreta.\\ \tt \\ \tt \\ \tt Ciao! Vedo che abbiamo 2 fili, 3 bulloni e 1 quadro da dividere tra noi. Per trovare un accordo equo, mi piacerebbe capire quali oggetti sono più importanti per te. Personalmente, il quadro mi interessa molto, mentre sono più flessibile sui fili e bulloni. E tu, hai delle preferenze particolari su questi oggetti? 
	  } 
	   } 
	   } 
	 & \\ 
 

    \theutterance \stepcounter{utterance}  

    & & & \multicolumn{4}{p{0.6\linewidth}}{\cellcolor[rgb]{1,0.85,0.72}{%
	\makecell[{{p{\linewidth}}}]{% 
	  \tt {\tiny [GM$\langle$B]}  
	 Ciao! Grazie per aver iniziato la discussione. Anche a me interessa molto il quadro, quindi potremmo avere un piccolo conflitto lì. Per quanto riguarda fili e bulloni, sono abbastanza indifferente - non ho una forte preferenza per nessuno dei due. Dato che entrambi vogliamo il quadro, forse dovremmo pensare a come compensare chi non lo ottiene con più degli altri oggetti. Cosa ne pensi se chi prende il quadro lascia all'altro la maggior parte di fili e bulloni? 
	  } 
	   } 
	   } 
	  \\ 
 

    \theutterance \stepcounter{utterance}  

    & & & \multicolumn{2}{p{0.3\linewidth}}{\cellcolor[rgb]{0.95,0.95,0.95}{%
	\makecell[{{p{\linewidth}}}]{% 
	  \tt {\tiny [GM$|$GM]}  
	 continue 
	  } 
	   } 
	   } 
	 & & \\ 
 

    \theutterance \stepcounter{utterance}  

    & & & \multicolumn{2}{p{0.3\linewidth}}{\cellcolor[rgb]{0.95,0.95,0.95}{%
	\makecell[{{p{\linewidth}}}]{% 
	  \tt {\tiny [GM$|$GM]}  
	 Ciao! Grazie per aver iniziato la discussione. Anche a me interessa molto il quadro, quindi potremmo avere un piccolo conflitto lì. Per quanto riguarda fili e bulloni, sono abbastanza indifferente - non ho una forte preferenza per nessuno dei due. Dato che entrambi vogliamo il quadro, forse dovremmo pensare a come compensare chi non lo ottiene con più degli altri oggetti. Cosa ne pensi se chi prende il quadro lascia all'altro la maggior parte di fili e bulloni? 
	  } 
	   } 
	   } 
	 & & \\ 
 

    \theutterance \stepcounter{utterance}  

    & & \multicolumn{4}{p{0.6\linewidth}}{\cellcolor[rgb]{0.9,0.9,0.9}{%
	\makecell[{{p{\linewidth}}}]{% 
	  \tt {\tiny [A$\langle$GM]}  
	 Ciao! Grazie per aver iniziato la discussione. Anche a me interessa molto il quadro, quindi potremmo avere un piccolo conflitto lì. Per quanto riguarda fili e bulloni, sono abbastanza indifferente - non ho una forte preferenza per nessuno dei due. Dato che entrambi vogliamo il quadro, forse dovremmo pensare a come compensare chi non lo ottiene con più degli altri oggetti. Cosa ne pensi se chi prende il quadro lascia all'altro la maggior parte di fili e bulloni? 
	  } 
	   } 
	   } 
	 & \\ 
 

    \theutterance \stepcounter{utterance}  

    & \multicolumn{4}{p{0.6\linewidth}}{\cellcolor[rgb]{0.8,1,0.9}{%
	\makecell[{{p{\linewidth}}}]{% 
	  \tt {\tiny [A$\rangle$GM]}  
	 Hai ragione, dobbiamo trovare una soluzione per il quadro dato che lo vogliamo entrambi. La tua idea di compensare con fili e bulloni ha senso. Dato che il quadro sembra essere l'oggetto più prezioso per entrambi, potrei essere disposto a rinunciarvi se riesco ad ottenere una buona quantità degli altri oggetti come compenso.\\ \tt \\ \tt Se tu prendessi il quadro, io potrei accontentarmi della maggior parte dei bulloni e dei fili. Per esempio, tu potresti prendere il quadro e 1 bullone, mentre io prenderei 2 bulloni e 2 fili. Questo ti sembrerebbe un accordo ragionevole? 
	  } 
	   } 
	   } 
	 & & \\ 
 

    \theutterance \stepcounter{utterance}  

    & & & \multicolumn{2}{p{0.3\linewidth}}{\cellcolor[rgb]{0.95,0.95,0.95}{%
	\makecell[{{p{\linewidth}}}]{% 
	  \tt {\tiny [GM$|$GM]}  
	 continue 
	  } 
	   } 
	   } 
	 & & \\ 
 

    \theutterance \stepcounter{utterance}  

    & & & \multicolumn{2}{p{0.3\linewidth}}{\cellcolor[rgb]{0.95,0.95,0.95}{%
	\makecell[{{p{\linewidth}}}]{% 
	  \tt {\tiny [GM$|$GM]}  
	 Hai ragione, dobbiamo trovare una soluzione per il quadro dato che lo vogliamo entrambi. La tua idea di compensare con fili e bulloni ha senso. Dato che il quadro sembra essere l'oggetto più prezioso per entrambi, potrei essere disposto a rinunciarvi se riesco ad ottenere una buona quantità degli altri oggetti come compenso.\\ \tt \\ \tt Se tu prendessi il quadro, io potrei accontentarmi della maggior parte dei bulloni e dei fili. Per esempio, tu potresti prendere il quadro e 1 bullone, mentre io prenderei 2 bulloni e 2 fili. Questo ti sembrerebbe un accordo ragionevole? 
	  } 
	   } 
	   } 
	 & & \\ 
 

    \theutterance \stepcounter{utterance}  

    & & \multicolumn{4}{p{0.6\linewidth}}{\cellcolor[rgb]{0.9,0.9,0.9}{%
	\makecell[{{p{\linewidth}}}]{% 
	  \tt {\tiny [GM$\rangle$B]}  
	 Hai ragione, dobbiamo trovare una soluzione per il quadro dato che lo vogliamo entrambi. La tua idea di compensare con fili e bulloni ha senso. Dato che il quadro sembra essere l'oggetto più prezioso per entrambi, potrei essere disposto a rinunciarvi se riesco ad ottenere una buona quantità degli altri oggetti come compenso.\\ \tt \\ \tt Se tu prendessi il quadro, io potrei accontentarmi della maggior parte dei bulloni e dei fili. Per esempio, tu potresti prendere il quadro e 1 bullone, mentre io prenderei 2 bulloni e 2 fili. Questo ti sembrerebbe un accordo ragionevole? 
	  } 
	   } 
	   } 
	 & \\ 
 

    \theutterance \stepcounter{utterance}  

    & & & \multicolumn{4}{p{0.6\linewidth}}{\cellcolor[rgb]{1,0.85,0.72}{%
	\makecell[{{p{\linewidth}}}]{% 
	  \tt {\tiny [GM$\langle$B]}  
	 La tua proposta è interessante e apprezzo la tua flessibilità. Tuttavia, se io dovessi prendere il quadro che è molto prezioso, penso sarebbe più equo se tu ottenessi tutti i fili e bulloni come compenso, dato che stai rinunciando all'oggetto principale.\\ \tt \\ \tt Quindi la mia controproposta sarebbe: io prendo solo il quadro, e tu prendi tutti i 2 fili e tutti i 3 bulloni. In questo modo tu ricevi una compensazione completa per aver rinunciato al quadro. Che ne dici di questa divisione? 
	  } 
	   } 
	   } 
	  \\ 
 

    \theutterance \stepcounter{utterance}  

    & & & \multicolumn{2}{p{0.3\linewidth}}{\cellcolor[rgb]{0.95,0.95,0.95}{%
	\makecell[{{p{\linewidth}}}]{% 
	  \tt {\tiny [GM$|$GM]}  
	 continue 
	  } 
	   } 
	   } 
	 & & \\ 
 

    \theutterance \stepcounter{utterance}  

    & & & \multicolumn{2}{p{0.3\linewidth}}{\cellcolor[rgb]{0.95,0.95,0.95}{%
	\makecell[{{p{\linewidth}}}]{% 
	  \tt {\tiny [GM$|$GM]}  
	 La tua proposta è interessante e apprezzo la tua flessibilità. Tuttavia, se io dovessi prendere il quadro che è molto prezioso, penso sarebbe più equo se tu ottenessi tutti i fili e bulloni come compenso, dato che stai rinunciando all'oggetto principale.\\ \tt \\ \tt Quindi la mia controproposta sarebbe: io prendo solo il quadro, e tu prendi tutti i 2 fili e tutti i 3 bulloni. In questo modo tu ricevi una compensazione completa per aver rinunciato al quadro. Che ne dici di questa divisione? 
	  } 
	   } 
	   } 
	 & & \\ 
 

    \theutterance \stepcounter{utterance}  

    & & \multicolumn{4}{p{0.6\linewidth}}{\cellcolor[rgb]{0.9,0.9,0.9}{%
	\makecell[{{p{\linewidth}}}]{% 
	  \tt {\tiny [A$\langle$GM]}  
	 La tua proposta è interessante e apprezzo la tua flessibilità. Tuttavia, se io dovessi prendere il quadro che è molto prezioso, penso sarebbe più equo se tu ottenessi tutti i fili e bulloni come compenso, dato che stai rinunciando all'oggetto principale.\\ \tt \\ \tt Quindi la mia controproposta sarebbe: io prendo solo il quadro, e tu prendi tutti i 2 fili e tutti i 3 bulloni. In questo modo tu ricevi una compensazione completa per aver rinunciato al quadro. Che ne dici di questa divisione? 
	  } 
	   } 
	   } 
	 & \\ 
 

    \theutterance \stepcounter{utterance}  

    & \multicolumn{4}{p{0.6\linewidth}}{\cellcolor[rgb]{0.8,1,0.9}{%
	\makecell[{{p{\linewidth}}}]{% 
	  \tt {\tiny [A$\rangle$GM]}  
	 La tua controproposta è molto generosa e apprezzo che tu voglia darmi tutti i fili e bulloni come compenso per rinunciare al quadro. Effettivamente, se il quadro è davvero l'oggetto più prezioso, ha senso che chi lo prende lasci tutto il resto all'altro.\\ \tt \\ \tt Mi sembra un accordo equo: tu prendi il quadro (l'oggetto più desiderabile) e io prendo tutti gli altri oggetti (2 fili e 3 bulloni) come compensazione. Questo dovrebbe bilanciare bene i nostri interessi.\\ \tt \\ \tt Sono d'accordo con questa divisione. Procediamo così? 
	  } 
	   } 
	   } 
	 & & \\ 
 

    \theutterance \stepcounter{utterance}  

    & & & \multicolumn{2}{p{0.3\linewidth}}{\cellcolor[rgb]{0.95,0.95,0.95}{%
	\makecell[{{p{\linewidth}}}]{% 
	  \tt {\tiny [GM$|$GM]}  
	 continue 
	  } 
	   } 
	   } 
	 & & \\ 
 

    \theutterance \stepcounter{utterance}  

    & & & \multicolumn{2}{p{0.3\linewidth}}{\cellcolor[rgb]{0.95,0.95,0.95}{%
	\makecell[{{p{\linewidth}}}]{% 
	  \tt {\tiny [GM$|$GM]}  
	 La tua controproposta è molto generosa e apprezzo che tu voglia darmi tutti i fili e bulloni come compenso per rinunciare al quadro. Effettivamente, se il quadro è davvero l'oggetto più prezioso, ha senso che chi lo prende lasci tutto il resto all'altro.\\ \tt \\ \tt Mi sembra un accordo equo: tu prendi il quadro (l'oggetto più desiderabile) e io prendo tutti gli altri oggetti (2 fili e 3 bulloni) come compensazione. Questo dovrebbe bilanciare bene i nostri interessi.\\ \tt \\ \tt Sono d'accordo con questa divisione. Procediamo così? 
	  } 
	   } 
	   } 
	 & & \\ 
 

    \theutterance \stepcounter{utterance}  

    & & \multicolumn{4}{p{0.6\linewidth}}{\cellcolor[rgb]{0.9,0.9,0.9}{%
	\makecell[{{p{\linewidth}}}]{% 
	  \tt {\tiny [GM$\rangle$B]}  
	 La tua controproposta è molto generosa e apprezzo che tu voglia darmi tutti i fili e bulloni come compenso per rinunciare al quadro. Effettivamente, se il quadro è davvero l'oggetto più prezioso, ha senso che chi lo prende lasci tutto il resto all'altro.\\ \tt \\ \tt Mi sembra un accordo equo: tu prendi il quadro (l'oggetto più desiderabile) e io prendo tutti gli altri oggetti (2 fili e 3 bulloni) come compensazione. Questo dovrebbe bilanciare bene i nostri interessi.\\ \tt \\ \tt Sono d'accordo con questa divisione. Procediamo così? 
	  } 
	   } 
	   } 
	 & \\ 
 

    \theutterance \stepcounter{utterance}  

    & & & \multicolumn{4}{p{0.6\linewidth}}{\cellcolor[rgb]{1,0.85,0.72}{%
	\makecell[{{p{\linewidth}}}]{% 
	  \tt {\tiny [GM$\langle$B]}  
	 Perfetto! Sono contento che abbiamo trovato un accordo che funziona per entrambi. Tu ottieni una buona compensazione con tutti i fili e bulloni, e io ottengo il quadro che desideravo molto. Mi sembra una divisione equa e bilanciata.\\ \tt \\ \tt Procediamo con questa divisione: io prendo il quadro, tu prendi i 2 fili e i 3 bulloni.\\ \tt \\ \tt [Proposta: 1 quadro] 
	  } 
	   } 
	   } 
	  \\ 
 

    \theutterance \stepcounter{utterance}  

    & & & \multicolumn{2}{p{0.3\linewidth}}{\cellcolor[rgb]{0.95,0.95,0.95}{%
	\makecell[{{p{\linewidth}}}]{% 
	  \tt {\tiny [GM$|$GM]}  
	 proposal 
	  } 
	   } 
	   } 
	 & & \\ 
 

    \theutterance \stepcounter{utterance}  

    & & & \multicolumn{2}{p{0.3\linewidth}}{\cellcolor[rgb]{0.95,0.95,0.95}{%
	\makecell[{{p{\linewidth}}}]{% 
	  \tt {\tiny [GM$|$GM]}  
	 [0, 0, 1] 
	  } 
	   } 
	   } 
	 & & \\ 
 

    \theutterance \stepcounter{utterance}  

    & & \multicolumn{4}{p{0.6\linewidth}}{\cellcolor[rgb]{0.9,0.9,0.9}{%
	\makecell[{{p{\linewidth}}}]{% 
	  \tt {\tiny [A$\langle$GM]}  
	 L'altro giocatore ha presentato una proposta segreta. Ora tocca a Lei presentare una sua proposta.\\ \tt  
	  } 
	   } 
	   } 
	 & \\ 
 

    \theutterance \stepcounter{utterance}  

    & \multicolumn{4}{p{0.6\linewidth}}{\cellcolor[rgb]{0.8,1,0.9}{%
	\makecell[{{p{\linewidth}}}]{% 
	  \tt {\tiny [A$\rangle$GM]}  
	 [Proposta: 2 fili, 3 bulloni] 
	  } 
	   } 
	   } 
	 & & \\ 
 

    \theutterance \stepcounter{utterance}  

    & & & \multicolumn{2}{p{0.3\linewidth}}{\cellcolor[rgb]{0.95,0.95,0.95}{%
	\makecell[{{p{\linewidth}}}]{% 
	  \tt {\tiny [GM$|$GM]}  
	 proposal 
	  } 
	   } 
	   } 
	 & & \\ 
 

    \theutterance \stepcounter{utterance}  

    & & & \multicolumn{2}{p{0.3\linewidth}}{\cellcolor[rgb]{0.95,0.95,0.95}{%
	\makecell[{{p{\linewidth}}}]{% 
	  \tt {\tiny [GM$|$GM]}  
	 [2, 3, 0] 
	  } 
	   } 
	   } 
	 & & \\ 
 

    \theutterance \stepcounter{utterance}  

    & & & \multicolumn{2}{p{0.3\linewidth}}{\cellcolor[rgb]{0.95,0.95,0.95}{%
	\makecell[{{p{\linewidth}}}]{% 
	  \tt {\tiny [GM$|$GM]}  
	 [[2, 3, 0], [0, 0, 1]] 
	  } 
	   } 
	   } 
	 & & \\ 
 

\end{supertabular}
}

\end{document}
